% ══════════════════════════════════════════════════════════════════════════════
\section{Empirical Strategy}
\label{sec:empirical}
% ══════════════════════════════════════════════════════════════════════════════

\subsection{Identification Challenge}

The causal identification of the effect of crop yields on conflict is complicated by a fundamental endogeneity problem: conflict itself destroys agricultural production. Armed conflict causes labour displacement, asset destruction, disruption of planting and harvest activities, and breakdown of market supply chains, all of which reduce measured crop yields. Studies that regress observed yields on conflict---or conflict on observed yields---therefore face a simultaneity bias whose direction is unclear a priori. Positive conflict $\rightarrow$ yield coefficients may be attenuated or reversed by the downward pressure of conflict on production, making OLS estimates with observed yields uninformative about the causal effect of interest.

Additionally, selection and omitted variable bias are potential concerns. Districts with historically fragile institutions, weak governance, and high ethnic heterogeneity may exhibit both lower average yields (due to insecurity, disinvestment, and limited access to markets and inputs) and higher average conflict. If these time-invariant characteristics are not controlled for, OLS estimates will confound the causal effect of yield shocks with the correlation driven by institutional quality. Similarly, common shocks---global commodity price movements, El Ni\~no events, pandemic-era disruptions---affect both yields and conflict simultaneously across all districts in a given year.

\subsection{Predicted Yield as the Regressor}

This paper's primary identification strategy exploits the distinction between \textit{observed} and \textit{predicted} crop yields. Observed yields, measured from household surveys, are endogenous to conflict: conflict reduces production and hence measured yields. Predicted yields, constructed from satellite and environmental inputs using a machine-learning model trained on survey observations, are not mechanically affected by conflict outcomes in the same way. Conflict affects whether a household can plant, cultivate, and harvest---and hence actual output---but it does not directly alter the satellite vegetation signal or the weather realised in a given growing season. The predicted yield for a district-year therefore reflects the environmental potential for crop production, stripped of the realised shortfalls induced by conflict and other idiosyncratic disruptions.

This approach is analogous to an intent-to-treat (ITT) analysis: the predicted yield captures the yield that \textit{would have been realised} given the environmental conditions, in the absence of any productivity-depressing shocks including conflict. Regressing forward conflict on predicted yield tests whether districts in environments conducive to higher agricultural output experience more or less subsequent conflict, conditional on fixed district characteristics and common time trends.

The exclusion restriction required for a causal interpretation is that satellite-derived vegetation signals and weather realisations in a district-year affect subsequent conflict only through their effect on crop productivity---not through independent channels. This assumption could be violated if, for example, temperature drives conflict directly through physiological pathways \citep{BurkeHsiangMiguel2015} or if vegetation greenness correlates with terrain features that independently predict conflict. To minimise this concern, the prediction models are evaluated using spatial out-of-fold cross-validation across distinct geographic clusters, ensuring that predictions reflect agronomically relevant variation rather than overfitting to local features correlated with conflict.

\subsection{Yield Prediction Models}

Yield predictions are generated using XGBoost models \citep{ChenGuestrin2016} trained on plot-level LSMS-ISA yield observations. A hybrid model selection approach is employed, choosing between country-specific and global models based on the best spatial out-of-fold $R^2$ for each country. The selected models are: country-specific models for Ethiopia ($R^2 = 0.066$), Malawi ($R^2 = 0.178$), Tanzania ($R^2 = 0.092$), and Uganda ($R^2 = 0.016$); and the global model for Mali ($R^2 = 0.192$) and Nigeria ($R^2 = 0.018$), where the global model outperforms the country-specific alternative.

Table~\ref{tab:model_r2} summarises model performance by country. The predictive accuracy of the yield models varies substantially. Malawi and Mali achieve the highest spatial $R^2$ values (0.178 and 0.192, respectively), while Uganda (0.016) and Nigeria (0.018) have near-zero predictive accuracy. This variation is attributable to differences in the complexity and heterogeneity of agricultural production systems---Nigeria's 464 districts span diverse agro-ecological zones with complex multi-season cropping systems that are difficult to capture with a single set of satellite features---as well as to differences in training sample size and geographic diversity across survey waves.

\begin{table}[H]
\centering
\caption{Yield Prediction Model Performance by Country}
\label{tab:model_r2}
\begin{threeparttable}
\begin{tabular}{lccc}
\toprule
Country & Model & Spatial $R^2$ & Training Obs. \\
\midrule
Ethiopia  & XGB (country) & 0.066 & --- \\
Malawi    & XGB (country) & 0.178 & --- \\
Mali      & XGB (global)  & 0.192 & --- \\
Nigeria   & XGB (global)  & 0.018 & --- \\
Tanzania  & XGB (country) & 0.092 & --- \\
Uganda    & XGB (country) & 0.016 & --- \\
\midrule
Niger     & Both models   & $< 0$ & 16 \\
\bottomrule
\end{tabular}
\begin{tablenotes}
\small
\item Spatial $R^2$ is computed using spatial out-of-fold cross-validation, where folds are defined by geographic clusters. A negative spatial $R^2$ indicates that the model performs worse than a baseline prediction of the overall mean. Niger is excluded from all analyses due to negative $R^2$ under both country-specific and global models.
\end{tablenotes}
\end{threeparttable}
\end{table}

All models predict yields in log scale; predictions are back-transformed using the inverse log transformation ($\exp(\hat{y}) - 1$) for descriptive purposes. The log-predicted yield, $\log\widehat{y}_{it}$, enters the regression directly without back-transformation.

\subsection{Baseline Specification}

The baseline empirical specification is a two-way fixed-effects (TWFE) OLS panel regression:
\begin{equation}
  \text{Conflict}_{i,t+h} = \beta \cdot \log\widehat{y}_{it} + \delta \cdot \log\text{Pop}_{it} + \alpha_i + \gamma_t + \varepsilon_{it}
  \label{eq:baseline}
\end{equation}
where $i$ indexes Admin-2 districts, $t$ indexes years, and $h$ denotes the length of the forward conflict window (3, 6, or 12 months). The key regressor $\log\widehat{y}_{it}$ is the natural log of the ML-predicted maize yield for district $i$ in year $t$. $\log\text{Pop}_{it}$ is the natural log of district population. $\alpha_i$ are Admin-2 fixed effects that absorb all time-invariant district characteristics including geography, ethnic composition, historical institutions, and long-run mean productivity. $\gamma_t$ are year fixed effects that absorb common shocks to conflict and yields including global commodity price cycles, El Ni\~no events, and pandemic-era disruptions. Standard errors are clustered at the Admin-2 level to allow for serial correlation in district-level conflict and yields.

The coefficient of interest is $\beta$: the within-district, within-year effect of log predicted yield on conflict count. A negative $\beta$ would be consistent with the opportunity-cost channel (higher yields reduce conflict), while a positive $\beta$ would be consistent with the predation or resource-contestation channel.

\subsection{Count Model (TWFE Poisson)}

Given that the outcome variable is a non-negative integer count with substantial zero-inflation (84.2\% zeros), OLS estimates of Equation~\ref{eq:baseline} may be inefficient and could produce negative predicted values. As an alternative, I estimate the same specification using Poisson quasi-maximum likelihood (QMLE):
\begin{equation}
  \mathbb{E}[\text{Conflict}_{i,t+h} \mid \log\widehat{y}_{it}, \alpha_i, \gamma_t] = \exp\!\left(\beta \cdot \log\widehat{y}_{it} + \delta \cdot \log\text{Pop}_{it} + \alpha_i + \gamma_t\right)
  \label{eq:poisson}
\end{equation}
\citet{SantosSilvaTenreyro2006} demonstrate that Poisson QMLE is consistent even when the conditional mean is misspecified, as long as the mean function is correctly specified, making it robust to overdispersion and zero-inflation. Fixed effects in the Poisson model are estimated by the Mundlak device \citep{fixest2020}, which drops districts with zero conflict outcomes in all periods (322 singleton Admin-2 units are dropped, leaving 604 districts and 9,060 observations in the Poisson sample). Standard errors are again clustered at the Admin-2 level.

\subsection{Log-Log Specification}

To obtain an elasticity estimate and handle the skewed distribution of conflict counts, I also estimate:
\begin{equation}
  \log(1 + \text{Conflict}_{i,t+h}) = \beta \cdot \log\widehat{y}_{it} + \delta \cdot \log\text{Pop}_{it} + \alpha_i + \gamma_t + \varepsilon_{it}
  \label{eq:loglog}
\end{equation}
The coefficient $\beta$ in this specification is interpretable as the elasticity of conflict with respect to predicted yield. The $\log(1 + \cdot)$ transformation retains the zero observations while compressing the right tail of the conflict distribution.

\subsection{Instrumental Variables (2SLS)}

As an additional robustness check---and to assess the plausibility of using predicted yield as an instrument for observed yield---I implement a 2SLS specification in which the endogenous variable is log observed yield (from LSMS surveys), instrumented by log predicted yield. The first and second stages are:
\begin{align}
  \log\hat{y}^{\text{obs}}_{it} &= \pi \cdot \log\widehat{y}_{it} + \delta \cdot \log\text{Pop}_{it} + \alpha_i + \gamma_t + u_{it} \quad \text{(first stage)} \\
  \text{Conflict}_{i,t+h} &= \beta_{\text{IV}} \cdot \widehat{\log y^{\text{obs}}}_{it} + \delta \cdot \log\text{Pop}_{it} + \alpha_i + \gamma_t + \varepsilon_{it} \quad \text{(second stage)}
\end{align}
This IV sample is restricted to district-years where LSMS observed yields are available, representing approximately 18\% of the full panel (2,528 observations across 883 districts). The first-stage F-statistic is reported as a diagnostic for instrument strength.

\subsection{Robustness Checks}

I conduct a series of robustness checks to assess the sensitivity of the main results:
\begin{enumerate}
  \item \textbf{Alternative conflict windows:} Repeat the TWFE OLS and Poisson specifications with 6-month and 12-month forward conflict windows.
  \item \textbf{Binary outcome (LPM):} Replace the count outcome with a binary indicator for any conflict in the 3-month window, estimated by OLS (linear probability model).
  \item \textbf{Fatality count (Poisson):} Replace conflict event count with total fatalities as the outcome, estimated by Poisson QMLE.
  \item \textbf{IV sample variation:} Estimate the 2SLS specification dropping Uganda (sparse IV coverage), dropping Uganda and Tanzania, and restricting to Ethiopia, Malawi, and Mali (the three countries with highest model $R^2$).
\end{enumerate}
